\documentclass[11pt]{article}
\usepackage[margin=1in]{geometry}
\usepackage{amsmath,amssymb,amsthm,mathtools,booktabs,longtable,array}
\usepackage[T1]{fontenc}
\usepackage{lmodern,microtype,mathrsfs}
\usepackage{hyperref}
\hypersetup{colorlinks=true,linkcolor=blue,citecolor=blue,urlcolor=blue}
\usepackage{listings}
\lstset{basicstyle=\ttfamily\small,breaklines=true,columns=fullflexible}
\newtheorem{theorem}{Theorem}[section]
\newtheorem{lemma}[theorem]{Lemma}
\newtheorem{proposition}[theorem]{Proposition}
\newtheorem{corollary}[theorem]{Corollary}
\theoremstyle{definition}\newtheorem{definition}[theorem]{Definition}
\theoremstyle{remark}\newtheorem{remark}[theorem]{Remark}
\newcommand{\K}{\mathcal K}
\newcommand{\HH}{\mathcal H}
\newcommand{\B}{\mathsf B}
\newcommand{\V}{\mathcal V}
\newcommand{\supp}{\operatorname{supp}}
\newcommand{\ord}{\operatorname{ord}}
\newcommand{\QR}{\operatorname{QR}}
\title{Pointwise angular forcing for marked Erd\H{o}s--Straus selectors\\
\large A continuation tranche: weighted exponent boxes, exact exceptions, and a finite covering certificate}
\author{Research transcript; no historical-priority claim}
\date{5 September 2026}
\begin{document}
\maketitle
\begin{abstract}
The proposed infinite-exterior/four-middle-shell dichotomy is not resolved here. We prove a pointwise angular-forcing theorem from actual bounded signed exponent boxes. Its optimal one-sided representative height gives explicit cofactor bounds, an exact marked construction, and complete inverse fibres. Seven concrete anchors reduce a possible counterexample to prescribed prime-factor residue sets; all finite angular exceptions are eliminated within the proposed hybrid. A balanced-factor family shows precisely why unweighted residue saturation does not itself imply chamber positivity. Independently verified modular certificates prove the proposed dichotomy for all $14\,215\,707$ hard primes at most $10^{10}$. The original finite exterior family with the four middle shells has exactly three failures in this enlarged range; each is repaired by an unbounded $(1,s)$ ray. No density theorem is used to infer any pointwise assertion.
\end{abstract}

\section{The exact question and its finite coordinates}
Let
\[
 C=\{1,121,169,289,361,529\},\qquad
 \HH=\{p\text{ prime}:p\bmod840\in C\}.
\]
Write
\[
 \alpha=\frac{\sqrt3-\sqrt2}{1+\sqrt2},\quad
 \beta=\frac{\sqrt3+\sqrt2}{1+\sqrt2},\quad
 \gamma=1+\sqrt2,\qquad
 \K=(0,\alpha)\cup(\beta,\gamma).
\]
Here $1/8<\alpha<1/7$, and $[3/2,2]\subset\K$.
The natural closure in this tranche means exactly
\begin{equation}\label{eq:closure}
 \V_\infty=
 \{(r,s)\in\mathbb Z_{>0}^2:\gcd(r,s)=1,\ 3s\le2r\le4s\}
 \ \cup\ \{(1,s):s\ge8\}.
\end{equation}
It does not add cross-mediants between the two disconnected portions of the chamber.

\begin{lemma}[Exact Farey domain]\label{lem:farey}
All primitive vectors in the first set in \eqref{eq:closure} are obtained by repeated insertion of mediants between $(2,1)$ and $(3,2)$.
\end{lemma}
\begin{proof}
The basis matrix with these columns has determinant one. A vector $(r,s)$ has basis coordinates
\[
 A=2r-3s,\qquad B=2s-r,
\quad (r,s)=A(2,1)+B(3,2).
\]
The interval inequalities are exactly $A,B\ge0$, and $\gcd(A,B)=1$. At an endpoint one coordinate is zero and the other one. If $A=B$, coprimality gives $A=B=1$, the mediant. If $A>B$, replace the current endpoint basis $(v,w)$ by $(v,v+w)$; the coordinates become $(A-B,B)$ and their sum decreases. If $B>A$, replace the basis by $(v+w,w)$, obtaining $(A,B-A)$. The determinant remains one. This subtractive Euclidean recursion terminates at an endpoint or a mediant, proving that the vector occurs in a finite refinement without assuming a common depth.
\end{proof}

The question is whether every $p\in\HH$ has an exterior state on $\V_\infty$, or a middle state with slope in $\K$ and $R\in\{11,19,23,39\}$. The marking is always
\[
 (p,R,a,u,h,r,s,\mathrm{channel},\kappa\text{ or }\lambda;
            x,y,z) ,
\]
where the last three coordinates retain their raw order. For $0<R<p$, $R\equiv3\pmod4$, set $a=(p+R)/4$. For each $u\mid a^2$, the unique normalization is
\begin{equation}\label{eq:normalize}
 d=\gcd(a,u),\quad r=u/d,\quad s=a/d,\quad h=d/r;
 \qquad a=hrs,\quad u=hr^2,\quad \gcd(r,s)=1.
\end{equation}
For a prime power $q^e\parallel a$, $q^f\parallel u$, $0\le f\le2e$, the exponents in $(h,r,s)$ are
\[
 (e-|f-e|,(f-e)_+,(e-f)_+).
\]
In particular no unavailable multiplicity is introduced. Also $u/a=r/s$.

The two exact channels are
\begin{align}
 E:&\quad R\mid4u+1,
 &\kappa&=(pr+s)/R,
 &(x,y,z)&=(hrs,hs\kappa,phr\kappa),\label{eq:E}\\
 M:&\quad R\mid4u+p,
 &\lambda&=(r+s)/R,
 &(x,y,z)&=(hrs,phs\lambda,phr\lambda).\label{eq:M}
\end{align}
Since $p\equiv4hrs\pmod R$ and $h,r,s$ are units modulo $R$, the divisibilities are respectively equivalent to $R\mid pr+s$ and $R\mid r+s$. Multiplication by $phrs\kappa$ or $phrs\lambda$ proves the reciprocal identity. These inherited coordinates and their literature crosswalk are isolated in Section~\ref{sec:antecedents}; the results below concern additional angular and pointwise information.

\section{Weighted angular forcing from an actual exponent box}
\begin{definition}
Let $R\ge3$ be odd, and let $H$ be an index-two subgroup of
$G=(\mathbb Z/R\mathbb Z)^\times$ with $-1\notin H$. An integer $A>0$ is a \emph{saturated anchor} for $(R,H)$ if $\gcd(A,R)=1$, all its prime factors lie in $H$, and its actual signed box
\begin{equation}\label{eq:box}
 \mathscr B_R(A)=
 \left\{nd^{-1}\pmod R:\ n,d>0,\ \gcd(n,d)=1,\ nd\mid A\right\}=H,
\end{equation}
Equivalently, if $A=\prod q^{e_q}$, the allowed exponent of $q$ in $n/d$ is each integer in $[-e_q,e_q]$. Define the one-sided representative height
\begin{equation}\label{eq:height}
 \B_R(A)=\max_{v\in H}\ 
 \min_{\substack{nd\mid A,\ \gcd(n,d)=1\\nd^{-1}\equiv v\ (R)}}\frac d n.
\end{equation}
Both extrema are over nonempty finite sets. In particular $\B_R(A)\le A$.
\end{definition}

\begin{theorem}[Pointwise weighted angular forcing]\label{thm:forcing}
Let $p>R$ satisfy $p\equiv1\pmod4$, $R\equiv3\pmod4$, and
$\gcd(a,R)=1$, where $a=(p+R)/4$. Primality is not needed for this construction. Let $A\mid a$ be a saturated anchor, put $b=a/A$, and let $0<\eta\le\alpha$. Set $T=\B_R(A)/\eta$.

If $b\bmod R\notin H$ and $b>T$, a middle state exists with $0<r/s<\eta$. The same conclusion holds if $b\bmod R\in H$, $b$ has an actual prime factor outside $H$, and $b>T^2$.
Consequently, if $a$ has a prime factor outside $H$, the single sufficient inequality
\begin{equation}\label{eq:coarseforce}
 a>A\max\{1,\B_R(A)/\eta\}^2
\end{equation}
forces such a state. In particular $a>64A^3$ suffices at $\eta=1/8$.
If all prime factors of $a$ lie in $H$, there is no middle state at that shell, irrespective of angle.
\end{theorem}
\begin{proof}
The quotient $G/H$ has order two, and its nonidentity coset is $-H$. We first find an actual divisor $t\mid b$, outside $H$, such that $t>T$.
If $b\notin H$, take $t=b$. Otherwise choose an actual prime $q\mid b$ outside $H$. If $q>T$, take $t=q$; if $q\le T$, take $t=b/q$. Then $b/q$ is outside $H$, and $b>T^2$ gives $b/q>T$. This is valid even for repeated prime factors: exactly one occurrence of $q$ has been removed.

The residue $-t^{-1}$ is in $H$. Choose a pair $(n,d)$ in the actual anchor box representing this residue with $d/n\le\B_R(A)$. Put
\begin{equation}\label{eq:packetmap}
 g=\gcd(d,tn),\quad r=d/g,\quad s=tn/g,\quad
 h=\frac{ag^2}{dtn},\quad
 u=\frac{ad}{tn},\quad
 \lambda=\frac{d+tn}{gR}.
\end{equation}
Because $nd\mid A$ and $t\mid b$, $dtn\mid a$. Thus $h$ and $u$ are positive integers, $a=hrs$, $u=hr^2$, and $\gcd(r,s)=1$. Also $R\mid d+tn$. Since $g\mid a$ and $\gcd(a,R)=1$, division by $g$ preserves the divisibility; hence $\lambda$ is integral. Formula~\eqref{eq:M} now gives the full ordered solution. Finally
\[
 \frac rs=\frac d{tn}\le\frac{\B_R(A)}t<\eta\le\alpha.
\]
For \eqref{eq:coarseforce}, if $T\ge1$ then $b>T^2$ also gives $b>T$; if $T<1$, every positive integer $b$ exceeds $T$. The bound $\B_R(A)\le A$ gives the last numerical assertion.

Conversely, a middle state would have $r,s\mid a$ and $r/s\equiv-1\pmod R$. If all prime factors of $a$ lie in $H$, both $r$ and $s$, with every available exponent, lie in $H$; their quotient cannot equal $-1$.
\end{proof}

\begin{corollary}[Necessary cofactor inequalities]\label{cor:heightobstruction}
Suppose the proposed hybrid fails at $p\in\HH$. Fix $R\in\{11,19,23,39\}$ and any saturated anchor $A\mid a_R=(p+R)/4$. If $a_R$ has a prime factor outside its anchor subgroup, then, with $B=\B_R(A)$,
\[
 a_R/A\notin H\ \Longrightarrow\ a_R\le8AB,
 \qquad
 a_R/A\in H\ \Longrightarrow\ a_R\le64AB^2.
\]
These inequalities apply to every such divisor $A$, not just to the seven fixed anchors below.
\end{corollary}
\begin{proof}
Apply the two branches of Theorem~\ref{thm:forcing} with $\eta=1/8$ and take contrapositives.
\end{proof}

\subsection{Domain, inverse fibres, and collision law}
Fix $p,R,A$ as above. The full packet domain for the rational threshold $1/8$ consists of
\[
 t\mid a/A,\quad nd\mid A,\quad\gcd(n,d)=1,\quad
 tn+d\equiv0\pmod R,\quad tn>8d.
\]
It maps by \eqref{eq:packetmap} into the marked middle states at $(p,R)$ with $r/s<1/8$. Not every such middle state need be in the image for a specified $A$.
For a target $(h,r,s)$ all preimages are given exactly by the anchor pairs $(n,d)$ for which
\begin{equation}\label{eq:packetinverse}
 t=\frac{ds}{nr}\in\mathbb Z_{>0},\qquad t\mid a/A,
 \qquad tn>8d,
\end{equation}
with the displayed congruence imposed (or inferred from the target). Then $r\mid d$, $g=d/r$, and \eqref{eq:packetmap} recovers the target. Indeed $r\mid ds=tnr$ and $\gcd(r,s)=1$ imply $r\mid d$, after which reduction of the ratio forces $g$.
Each $(n,d)$ permits at most one $t$, so every fibre is finite and has size at most $\prod_{q^{e_q}\parallel A}(2e_q+1)$.
Two packets collide precisely when
\begin{equation}\label{eq:collision}
 tnd'=t'n'd.
\end{equation}
This is equality of their reduced slopes, equivalently equality of their $u$ coordinates. Thus dropping packet coordinates loses a finite fibre; it does not erase the exterior/middle tag, shell, normalization, or denominator orientation.

\section{Seven explicit anchors and their exact angular exceptions}\label{sec:anchors}
For $R=11,19,23$, let $H_R$ be the quadratic residues. Put
\[
 H_{39}=\langle2\rangle=\{1,2,4,5,8,10,11,16,20,22,25,32\}.
\]
All these are index-two subgroups excluding $-1$.

\begin{theorem}[Exact hard-prime anchor reduction]\label{thm:seven}
The following table gives saturated anchors and their exact heights. In the last column are all $p\in\HH$ for which $A\mid a_R$, some prime factor of $a_R$ is outside $H_R$, yet no middle state at $R$ has slope in $\K$.
\begin{center}
\begin{tabular}{rrrrr}
\toprule
$R$&$A$&$\B_R(A)$&finite cutoff $256A\B_R(A)^2-R$&exceptions\\
\midrule
11&9&9&186613&none\\
23&12&3&27625&none\\
23&18&6&165865&1201\\
19&35&35&10975981&837601\\
39&40&$5/2$&63961&none\\
39&50&10&1279961&3361\\
39&64&32&16777177&none\\
\bottomrule
\end{tabular}
\end{center}
Every exception has a positive witness inside the proposed hybrid. Consequently any counterexample satisfies, for every row,
\begin{equation}\label{eq:puresupport}
 A\mid a_R\quad\Longrightarrow\quad
  q\bmod R\in H_R\quad\text{for every prime }q\mid a_R.
\end{equation}
\end{theorem}
\begin{proof}
The box computations are small and exact. Modulo $11$, $3$ has order $5$; the exponents $-2,\ldots,2$ fill $H_{11}$. Modulo $23$, $2$ has order $11$ and $3=2^8$; the rectangles
\[
 -2\le i\le2,\ -1\le j\le1,
 \qquad\text{and}\qquad
 -1\le i\le1,\ -2\le j\le2
\]
in the exponent $i+8j$ each cover $\mathbb Z/11\mathbb Z$. Modulo $19$, $5$ has order $9$ and $7=5^6$; $i+6j$, $i,j\in\{-1,0,1\}$, covers all nine classes. Modulo $39$, $2$ has order $12$ and $5=2^9$. The rectangles $|i|\le3,|j|\le1$ and $|i|\le1,|j|\le2$ in $i+9j$, and the interval $|i|\le6$, cover all twelve classes.

For the heights one enumerates those same boxes, taking the least $d/n$ in each residue. The accompanying JSON contains every pair, all signed exponents, and all minimizing representatives. This is exhaustive finite minimization, not subgroup generation. For example the minimizing pairs for $R=23,A=12$, in increasing order of $H_{23}$, are
\[
 (n,d)=(1,1),(2,1),(3,1),(4,1),(6,1),(1,3),
 (4,3),(12,1),(3,2),(2,3),(3,4).
\]
Their maximum $d/n$ is $3$, and the residue $8$ requires $(1,3)$, proving equality.

Theorem~\ref{thm:forcing} excludes all exceptions with
$a_R>64A\B_R(A)^2$. The remaining finite calculation runs
\begin{equation}\label{eq:finiteanchorloop}
 1\le b\le64\B_R(A)^2,\qquad p=4Ab-R.
\end{equation}
For each hard prime in this interval, factor $a_R=Ab$, test its actual prime factors against $H_R$, then enumerate every exponent $0\le f_q\le2v_q(a_R)$ for $u\mid a_R^2$. Test $R\mid4u+p$ and the strict integer positivity test in Appendix~\ref{app:positivity}. The exact counts are:
\begin{center}
\begin{tabular}{rrrrr}
\toprule
$R$&$A$&hard primes&with a non-$H_R$ factor&angular exceptions\\
\midrule
11&9&152&91&0\\
23&12&37&24&0\\
23&18&143&90&1\\
19&35&3754&2421&1\\
39&40&19&9&0\\
39&50&302&183&1\\
39&64&1061&682&0\\
\bottomrule
\end{tabular}
\end{center}
These are counts of primes, not counts of states. The verifier uses trial division for primality and complete divisor boxes, and regenerates this entire table. Thus the finite part is a deterministic computer-assisted lemma with the explicit bound \eqref{eq:finiteanchorloop}.

At the three exceptional shells, all oriented middle states are exactly the following; each entry lists $(u,h,r,s,\lambda)$:
\begin{align*}
 p=1201,\,R=23:&\quad (108,3,6,17,1),\ (867,3,17,6,1);\\
 p=837601,\,R=19:&\quad (37975,31,35,193,12),\ (47285,965,7,31,2),\\
 &\quad (927365,965,31,7,2),\ (1154719,31,193,35,12);\\
 p=3361,\,R=39:&\quad (125,5,5,34,1),\ (5780,5,34,5,1).
\end{align*}
Their repairs, with full coordinates in the JSON, are
\begin{center}
\begin{tabular}{rrcrrrrr}
\toprule
$p$&$R$&tag&$u$&$h$&$r$&$s$&$\lambda$ or $\kappa$\\
\midrule
1201&39&M&2&2&1&155&4\\
837601&11&M&81&1&9&23267&2116\\
3361&3&E&29&29&1&29&1130\\
\bottomrule
\end{tabular}
\end{center}
The raw ordered denominators are respectively
\[
 (310,1489240,9608),\quad
 (209403,41237586580172,15951273444),\quad
 (841,950330,110139970).
\]
Direct substitution or \eqref{eq:E}--\eqref{eq:M} verifies each. None is a counterexample to the proposed hybrid, proving \eqref{eq:puresupport}.
\end{proof}

\subsection{Hand classification for the first three anchors}
The finite classification at $R=11,23$ also has a small direct proof independent of the larger cutoff enumeration.

\begin{proposition}\label{prop:hand}
Among odd integers $a$ with $9\mid a$ and $\gcd(a,11)=1$ which contain a nonresidue prime factor, failure of chamber-positive middle selection at $R=11$ occurs exactly at
\[
 a\in\{171,513,3249,9747\}.
\]
Among positive integers $a$ coprime to $23$ containing a nonresidue prime factor, the analogous exceptional sets are $\{60,300\}$ when $12\mid a$, and $\{306,5202\}$ when $18\mid a$.
\end{proposition}
\begin{proof}
For an anchor $A$ and an outside prime $q>8A$, Theorem~\ref{thm:forcing}'s construction with $t=q$ gives a positive state whenever $Aq\mid a$. For the remaining primes $q\le8A$, inspect only the anchor box and the two orientations of $qn/d$. The exact bad-prime lists are
\[
 (R,A)=(11,9):\ \{2,19\};\qquad
 (23,12):\ \{5\};\qquad
 (23,18):\ \{17\}.
\]
The small-prime table in the certificate lists all positive packets and proves these finite assertions by integer inequalities.

For $(11,9)$, oddness excludes $2$, so a remaining bad $a$ can have no nonresidue prime other than $19$. If another QR prime $\ell\ne3$ divides $a$, then $t=19\ell>72$ is an available outside divisor of $a/9$. An exponent at least three on $19$ similarly supplies $t=19^3>72$. If $v_3(a)\ge4$, the actual divisor $3^4\cdot19=1539\equiv-1\pmod{11}$ gives $(r,s)=(1,1539)$. Thus only $a=3^e19^f$, $e=2,3$, $f=1,2$, remain.

For $(23,12)$, only the nonresidue $5$ can remain. Any further QR prime is $13$ or at least $29$, apart from $2,3$. The factor $13$ gives the actual divisor $390\equiv-1\pmod{23}$; a larger QR prime $\ell$ gives $t=5\ell>96$. The cases $v_2(a)\ge3$, $v_3(a)\ge2$, and $v_5(a)\ge3$ give respectively
\[
 (r,s)=(15,8),\quad(1,45),\quad(3,250),
\]
whose products divide $a$, sums are divisible by $23$, and slopes are in $\K$. Only $a=60,300$ remain.

For $(23,18)$ only the nonresidue $17$ can remain. Any other QR prime $\ell\ne2,3$ is at least $13$, so $t=17\ell>144$. The cases $v_2(a)\ge2$, $v_3(a)\ge3$, and $v_{17}(a)\ge3$ give the divisors
\[
 68,\qquad459,\qquad18\cdot17^3=88434,
\]
all $-1$ modulo $23$, and hence the positive rays $(1,s)$ in the middle channel. Only $a=306,5202$ remain.

For each listed $a$, complete divisor enumeration proves nonpositivity; these tiny boxes are also in the JSON. The corresponding $4a-R$ values in the $R=11$ case are $673,2041,12985,38977$, none a hard prime ($2041=13\cdot157$, $12985$ is divisible by $5$, and the other two are not hard residues). For $12\mid a$, $217=7\cdot31$ and $1177=11\cdot107$ are composite. For $18\mid a$, $4\cdot306-23=1201$ and $4\cdot5202-23=20785$ is divisible by $5$.
\end{proof}

\section{Additional pointwise restrictions and a genuine angular barrier}
\begin{proposition}[The shell-$3$ exterior dichotomy]\label{prop:R3}
For a hard prime $p$, put $a_3=(p+3)/4$. There is an exterior state at $R=3$ with $r=1,s\ge8$ if and only if $a_3$ has a prime factor $2\pmod3$.
\end{proposition}
\begin{proof}
We have $a_3\equiv1\pmod6$. If $q\equiv2\pmod3$ is its least such prime factor, the multiplicity of outside factors in $a_3$ is even. Hence $s=a_3/q\ge q$ and $s\equiv5\pmod6$. If $s<8$, then $q=s=5$, giving $p=97$, which is not hard. Thus $s\ge8$, and $(R,h,r,s)=(3,q,1,a_3/q)$ satisfies $R\mid4h+1$, giving the claimed exterior state. Conversely such a state has $h\equiv2\pmod3$, so an actual prime factor of $h\mid a_3$ is $2\pmod3$.
\end{proof}

\begin{corollary}[Explicit counterexample locus]\label{cor:locus}
Any counterexample to the proposed hybrid has all prime factors of $(p+3)/4$ congruent to $1\pmod3$, satisfies all seven implications \eqref{eq:puresupport}, and satisfies every variable-anchor inequality in Corollary~\ref{cor:heightobstruction}.
Moreover, unless $p\equiv1\pmod{144}$, at least one of the first three fixed anchor implications applies. On the entire hard class $p\equiv121\pmod{840}$, the $R=19,A=35$ implication applies.
\end{corollary}
\begin{proof}
Only the last assertions require calculation. Let $b=(p+23)/24$. Then
\[
 a_{11}=3(2b-1),\qquad a_{23}=6b.
\]
The three anchor conditions are respectively $b\equiv2\pmod3$, $2\mid b$, and $3\mid b$. None holds precisely when $b\equiv1\pmod6$, equivalently $p\equiv1\pmod{144}$. Finally $p\equiv121\pmod{840}$ implies $p+19\equiv0\pmod{140}$, hence $35\mid a_{19}$.
\end{proof}

The height obstruction also allows anchors not in the fixed list. For example, if $\ell>3$ is a prime in a nonidentity QR class modulo $11$, then $A=3\ell$ is saturated. Direct inspection gives
\begin{equation}\label{eq:variableheight}
\begin{array}{c|rrrr}
 \ell\bmod11&3&4&5&9\\\hline
 \B_{11}(3\ell)&3\ell&\ell/3&1&3.
\end{array}
\end{equation}
To verify this, use $3$ as generator of $H_{11}$: the second prime has exponent $1,-1,-2,2$ respectively, and inspect the nine exponents $i+cj$, $i,j=-1,0,1$. For completeness, in the residue order $1,3,4,5,9$ of $H_{11}$, the exact minima of $d/n$ are
\[
\begin{array}{c|ccccc}
\ell\bmod11&1&3&4&5&9\\\hline
3&3/\ell&1/\ell&3&3\ell&1/(3\ell)\\
4&1/(3\ell)&1/3&1/\ell&3/\ell&\ell/3\\
5&1&1/3&1/(3\ell)&1/\ell&3/\ell\\
9&1&3/\ell&3&1/(3\ell)&1/\ell
\end{array}
\]
Every entry is among the nine available monomials $3^{-i}\ell^{-j}$; comparison of the alternatives proves minimality. In the last row $\ell>9$, since $\ell>3$ is a prime congruent to $9$ modulo $11$. All other comparisons use $\ell>3$. Taking the row maxima proves \eqref{eq:variableheight}. Thus a nonresidue factor in $a_{11}$, combined with a QR prime $\ell\equiv5\pmod{11}$, forces a positive middle state as soon as $a_{11}>192\ell$; for $\ell\equiv9\pmod{11}$ the bound is $a_{11}>1728\ell$. These are linear, not cubic, bounds. In the outside-cofactor branch they improve further to $24\ell$ and $72\ell$, respectively.

\begin{theorem}[A balanced-factor obstruction to angular propagation]\label{thm:barrier}
Let $\ell,q$ be distinct primes other than $3,11$, satisfying
\[
 \ell\equiv3\pmod{11},\qquad q\equiv2\pmod{11},
 \qquad9\ell\le q\le21\ell.
\]
Put $a=3\ell q$ and $p=4a-11$. Whenever $p>11$ is prime, its complete middle fibre at $R=11$ consists of exactly the two slopes
\[
 \frac{3\ell}{q},\qquad\frac q{3\ell}.
\]
Neither is in $\K$. The signed QR box from $3\ell$ nevertheless fills the entire QR subgroup, and $q$ is a nonresidue.
\end{theorem}
\begin{proof}
Use primitive root $2$ modulo $11$. The exponents of $3,\ell,q$ are $8,8,1$ modulo $10$. In the actual box each exponent multiplier is $-1,0,1$. The equation
\[
 8(i+j)+k\equiv5\pmod{10},\qquad i,j,k\in\{-1,0,1\},
\]
has exactly the solutions $(1,1,-1)$ and $(-1,-1,1)$: inspect $i+j=-2,-1,0,1,2$; only the extreme sums work, with $k=1,-1$ respectively. Thus the only coprime divisor ratios equal to $-1$ modulo $11$ are $3\ell/q$ and its inverse, and $h=1$.
The first ratio lies in $[1/7,1/3]$, above $\alpha$ but below $\beta$; the inverse lies in $[3,7]$, above $\gamma$. Both fail the strict chamber inequalities. Saturation of $3\ell$ follows from $i+j\in\{-2,-1,0,1,2\}$ modulo $5$.
\end{proof}

A concrete hard prime is
\[
 p=2016361=12\cdot113\cdot1487-11,
 \qquad a=504093=3\cdot113\cdot1487.
\]
Its full $R=11$ fibre has
\[
 (u,h,r,s,\lambda)=(114921,1,339,1487,166),
 \quad(2211169,1,1487,339,166),
\]
with raw denominators
\[
 (504093,497722581962,113468698914),
 \quad(504093,113468698914,497722581962).
\]
Primality, the complete exponent box, and nonpositivity are deterministic certificate checks. There is no assertion that infinitely many prime values of the displayed bilinear expression exist. The family proves a conditional-in-parameters obstruction, not a hypothetical repair of the global selector.

\section{A certificate through $10^{10}$, independent of the discovery search}\label{sec:cover}
Let
\[
 \V_* = F_4\cup\{(1,s):8\le s\le128\}.
\]
Here $F_4$ has seventeen rays, so $|\V_*|=138$.

\begin{theorem}[Deterministic finite-range coverage]\label{thm:range}
There are exactly $14\,215\,707$ primes $p\le10^{10}$ in $\HH$. The exterior family $\V_*$ together with the four positive middle shells has exactly three failures in that interval:
\[
 11259889,\quad788653441,\quad1350237001.
\]
All three have exterior witnesses on $\V_\infty$. Consequently the proposed infinite-closure/four-shell dichotomy holds for every hard prime at most $10^{10}$.
\end{theorem}
\begin{proof}[Certificate and proof of its verification rule]
The supplied compressed JSON has $78\,332$ exterior rules and $414\,902$ middle rules. Its generator is not trusted: the standard-library verifier independently checks the validity of every rule and exhausts the prime universe by an exact sieve.

An exterior rule fixes primitive $(r,s)\in\V_\infty$ and $R\equiv3\pmod4$ with $\gcd(4r^2s,R)=1$. Put $m=4rs$ and
\begin{equation}\label{eq:Erule}
 t\equiv-s(mr)^{-1}\pmod R,\qquad
 p\equiv-R+mt\pmod{mR},\qquad p>R.
\end{equation}
Every such prime has $h=(p+R)/m$, $u=hr^2$ and $\kappa=(pr+s)/R$. Formula~\eqref{eq:E} reconstructs all marks. In the certificate, every rule belongs either to $\V_*$ or to one of the three added rays below.

A middle rule fixes $(R,u,d)$ with
\begin{equation}\label{eq:Mrule}
 R\in\{11,19,23,39\},\quad u\mid d^2,\quad\gcd(d,R)=1,
 \quad a\equiv0\pmod d,\quad a\equiv-u\pmod R.
\end{equation}
These congruences define one class $p=4a-R$ modulo $4dR$. If $u/a$ is in the lower chamber at some threshold $p_0>R$, it stays in the lower chamber for every larger $p$ in that class. If $u/a$ is in the upper chamber at both endpoints of a closed interval, it stays there throughout the interval, since $a$ increases with $p$ and $u/a$ decreases. Thus each rule stores a lower threshold, an upper-chamber interval, or both. The verifier checks the endpoints using strict integer inequalities. The conditions imply $u\mid a^2$, $\gcd(a,R)=1$, and $R\mid4u+p$; normalization \eqref{eq:normalize} reconstructs the full middle marking at every covered prime. This compression loses the discovery witness, but not the ability to recover a canonical complete ordered witness.

For completeness, the prime sieve has one flag for each integer $840k+c$, $c\in C$, up to the bound. For each prime $q\le10^5$ not dividing $840$, strike the unique $k$-class
\[
 k\equiv-c\,840^{-1}\pmod q
\]
starting at values $840k+c\ge q^2$. No member of these hard classes is divisible by $2,3,5,7$. The integer $1$ is explicitly deleted. Every composite in the universe has a prime divisor at most $10^5$; hence this is a complete deterministic sieve. Counting its surviving flags gives the stated prime count.

Each verified witness progression is intersected with each of the six hard progressions by the generalized Chinese remainder theorem. To make the counting proof explicit, $840k+c\equiv b\pmod M$ is soluble precisely when $g=\gcd(840,M)$ divides $b-c$. If it does, $k$ occupies the single class
\[
 k\equiv\frac{b-c}{g}(840/g)^{-1}\pmod{M/g}.
\]
The verifier clears exactly these flags inside the stored interval. After all base-family exterior and all middle rules, exactly the three displayed flags remain. Exact factorization of $rp+s$ for all $138$ base rays, together with complete middle divisor boxes at the four shells, proves that these three are genuine failures, rather than merely gaps in the compressed cover.

The additional exterior states are
\begin{center}
\begin{tabular}{rrrrrrr}
\toprule
$p$&$R$&$u=h$&$r$&$s$&$\kappa$&$a$\\
\midrule
11259889&151&3586&1&785&74574&2815010\\
788653441&7&698&1&282469&112705130&197163362\\
1350237001&3&983&1&343397&450193466&337559251\\
\bottomrule
\end{tabular}
\end{center}
All have $s\ge8$ and $0<R<p$. Their raw denominators are
\begin{align*}
 &(2815010,209926555740,3011146134757596),\\
 &(197163362,22221322345447060,62041931437810926340),\\
 &(337559251,151966969188053966,597534121519807663078).
\end{align*}
Their three rules clear the remaining flags. This proves the finite assertion, not a limiting claim.
\end{proof}

In fact all four middle fibres at each of these three primes are empty, even before imposing positivity. The JSON retains their factorizations, all $138$ base-ray factorizations, and every actual middle state (the lists are empty). State counts and prime counts are not interchanged.

\section{Complete termination for one candidate, and its limitation}
A future candidate can be tested without any Farey-depth cutoff. Since $u/a=r/s$, membership in the closure is exactly
\begin{equation}\label{eq:uclosure}
 3a\le2u\le4a
 \quad\text{or}\quad
 u\mid a\ \text{ and }\ 8u\le a.
\end{equation}
At each $R=3,7,\ldots<p$, the exterior congruence is
\[
 u\equiv u_0=\frac{3R-1}{4}\pmod R.
\]
In the upper branch enumerate these $u$ between $\lceil3a/2\rceil$ and $2a$ and test $u\mid a^2$; in the lower branch enumerate them up to $a/8$ and test $u\mid a$. In the lower branch $4u+1\ge3R$ and $8u\le a=(p+R)/4$ imply $23R\le p+8$, so no larger lower shell can contribute.

\begin{proposition}[Exact exterior hyperbola map]\label{prop:hyperbola}
There is a bijection between exterior markings $(p,R,a,u)$ and the tuples $(p,k,D,u)$ satisfying
\[
 k>0,\quad k\equiv3\pmod4,\quad D=(kp+1)/4,\quad
 u\mid D^2,\quad k\mid u+D,
 \quad0<\frac{4u+1}{k}<p.
\]
The maps are
\[
 k=\frac{4u+1}{R},\quad D=ka-u,
 \qquad
 R=\frac{4u+1}{k},\quad a=\frac{u+D}{k}.
\]
All fibres have size one when $p,u$ are retained. The normalization and channel tag remain unchanged.
\end{proposition}
\begin{proof}
Forwards, $4D=kp+1$ follows from $4a=p+R$. Also $\gcd(k,u)=1$, since any common divisor divides $4u+1=kR$ and $u$. Modulo $u$, $D\equiv ka$, whence $u\mid a^2$ implies $u\mid D^2$.
Conversely $4D-kp=1$ gives $\gcd(k,D)=1$; from $u\equiv-D\pmod k$ follows $\gcd(k,u)=1$. The same congruence gives $k\mid4u+1$. Then $4a=p+R$, $R\equiv3\pmod4$, and $u\mid D^2$ with $D\equiv ka\pmod u$ yields $u\mid a^2$. All maps are now forced, proving the inverse assertions.
\end{proof}

For the closure $u\le2a\le p-1$, so $kR=4u+1\le4p-3$. Choose any integer cutoff $0\le T<p$. Enumerate small shells $R\le T$ by divisors of $a^2$; enumerate the remaining states by
\[
 k\equiv3\pmod4,\qquad
 k\le\left\lfloor\frac{4p-3}{T+1}\right\rfloor,
 \qquad u\mid\left(\frac{kp+1}{4}\right)^2,
\]
retaining only $R>T$ and \eqref{eq:uclosure}. This is a disjoint exhaustive two-pass algorithm. Every divisor enumeration uses the exact exponent interval $0\le f_q\le2e_q$. The verifier implements both algorithms and compares their full outputs at small hard primes. No optimal-complexity claim is made; Elsholtz--Tao already discuss a different Type-I search algorithm in their Section~3.

Termination for testing one specified $p$ is not a terminating existence proof for every $p$. No assertion that the two-pass output is always nonempty has been inserted.

\section{Where the pointwise argument stops}\label{sec:stop}
The counterexample locus in Corollary~\ref{cor:locus} is not empty as a set of primes. For example $p=11259889$ satisfies every stated fixed-anchor support condition and the shell-$3$ support condition:
\[
 a_3=7\cdot402139,
 \qquad a_{11}=3^2\cdot5^2\cdot12511.
\]
Both factors of $a_3$ are $1\pmod3$; all factors of $a_{11}$ are QR modulo $11$. Among the seven fixed anchors only $(R,A)=(11,9)$ applies. Its other shifted factorizations are
\[
 a_{19}=11\cdot255907,\quad
 a_{23}=2\cdot3\cdot41\cdot11443,\quad
 a_{39}=2\cdot1407491.
\]
All four middle fibres are empty, as certified, so the variable-anchor necessary inequalities hold as well by Theorem~\ref{thm:forcing}. Thus this prime does not merely defeat a particular angular construction: these four middle shells cannot help at all. It is repaired by the exterior ray $(1,785)$ and shell $151$.

The next pointwise theorem therefore must force an exterior witness in the closure on the remaining weighted/support locus, or supply a logically complete alternative using one of the existing four middle shells when its fibre is nonempty. A subgroup membership argument cannot replace an actual bounded-exponent divisor; unweighted signed-box saturation cannot replace the angle, by Theorem~\ref{thm:barrier}. The current results prove neither that the remaining locus is finite nor that its exterior selector always terminates successfully. They locate any counterexample above $10^{10}$ and impose explicit additional factorization inequalities at every saturated anchor.

\section{Content-level antecedent crosswalk}\label{sec:antecedents}
The following distinguishes inherited identities from the angular continuation. Page numbers for Elsholtz--Tao refer to the 55-page arXiv manuscript, so that equation numbers give a version-stable second locator.

\paragraph{Yamamoto \cite{Yamamoto}.}
Lemma~2, equation~(4), printed p.~38, is the exterior congruence pair, with Yamamoto's $(a,b,q)=(r,s,R)$. The proof of Lemma~1 on the same page separates one and two $p$-divisible denominators. These facts are inherited, not new.

\paragraph{Monks--Velingker \cite{MV}.}
Theorem~2.1(c)--(e), pp.~2--3, supplies the divisibility/valuation classification and orientation context. Theorem~2.2 and Corollary~2.3, p.~4, with the proof of Theorem~2.2 on pp.~5--6, supply the two-unit-fraction divisor criterion. Its reduced-fraction form is the one used here: $\gcd(R,pa)=1$ and coprime divisor pairs. No inference from an unreduced numerator/denominator is needed.

\paragraph{Elsholtz--Tao \cite{ET}.}
For the exterior state, their Type-I coordinates in equations~(2.1)--(2.9), p.~11, are exactly
\[
 (a_{\rm ET},b_{\rm ET},c_{\rm ET},d_{\rm ET},e_{\rm ET},f_{\rm ET})
   =(r,\kappa,s,h,k,R),\quad k=(4hr^2+1)/R.
\]
Their map $\pi^{\rm I}$ gives $(z,x,y)$ in our order. The equation $ce=a+b$ is $sk=r+\kappa$, obtained by comparing $skR=s(4hr^2+1)$ with $Rr+pr+s$. Their Type-II coordinates in equations~(2.13)--(2.22), p.~14, are
\[
 (a_{\rm ET},b_{\rm ET},c_{\rm ET},d_{\rm ET},e_{\rm ET},f_{\rm ET})
  =(s,r,\lambda,h,R,4hs\lambda-1).
\]
Their map $\pi^{\rm II}$ gives exactly $(x,y,z)$ in our raw order. Propositions~2.2 and 2.6 are the corresponding coordinate descriptions. Thus neither the channel normalization nor its two-fraction identity is being assigned new priority.

\paragraph{Dahan \cite{Dahan}.}
Lemma~2.2 in Section~2.2, and Theorem~2.6/Corollary~2.7 in Section~2.3 of v1, give the reduced binary criterion and untwisted/twisted divisor branches. In the untwisted branch take $(n,c,K,d_1,d_2)=(p,R,a,r,s)$; in the twisted branch use the same tuple with congruence $pr+s\equiv0\pmod R$. Lemma~4.2 and Theorem~4.3 in Section~4.1 concern residue-support restrictions and sieve dimension, not the pointwise angular conclusion proved here. The repairs to the inherited sieve presentation are recorded separately in Appendix~\ref{app:sieve}.

\paragraph{Content search beyond the four requested antecedents.}
A search located the repository notes \cite{Repo11,Repo23} \emph{Exact $q=11$ filter}, Sections~3, 7--8, and \emph{Exact $q=23$ filter}, Sections~3--7, dated 15 August 2026 in the repository \texttt{chasebryan/centl}. They already discuss unweighted signed-box QR saturation, including $9\mid a_{11}$ and extra $2$- or $3$-valuation in $a_{23}$. The residue-saturation idea is therefore explicitly not claimed as new. Those notes concern the existence of a middle target without the chamber restriction. The present additions are the one-sided height, its actual-divisor cofactor construction and inverse fibres, exact angular exceptions, the balanced-factor angular obstruction, and the independent $10^{10}$ certificate. No exhaustive historical-priority assertion is made for these additions either. The proof of the angular results does not require accepting a repository status label or an unproved global claim.

\appendix
\section{Exact positivity test}\label{app:positivity}
For positive integers $r,s$, set
\[
 D=s^2-4rs-3r^2,\qquad E=3r^2-4rs-s^2.
\]
Then
\begin{align*}
 r/s<\alpha
 &\iff D>0\ \text{and}\ D^2>8r^2(r+s)^2,\\
 r/s<\gamma
 &\iff r\le s\ \text{or}\ (r-s)^2<2s^2,\\
 r/s>\beta
 &\iff r>s\ \text{and}\bigl(E\ge0\ \text{or}\
                    8r^2(r-s)^2>E^2\bigr).
\end{align*}
For the first equivalence rewrite $r/s+(r/s+1)\sqrt2<\sqrt3$, square positive sides, and isolate the remaining positive radical before squaring again. For the last, first $r/s>\beta>1$, then rewrite $r/s+(r/s-1)\sqrt2>\sqrt3$ and square positive sides. If $E\ge0$, the radical term gives strict positivity; otherwise the second square is equivalent. These sign conditions prevent all extraneous roots. The program uses these strict integer tests, never floating-point approximations.

\section{The two inherited sieve presentation repairs}\label{app:sieve}
No new density statement is used in this tranche. To specify the corrected notation in the inherited proof, fix once and for all its allowed-prime set $\mathcal P$ and write
\[
 \mathcal D=\{d\ge1:d\text{ squarefree and every prime divisor of }d
                      \text{ belongs to }\mathcal P\}.
\]
Every prime product is restricted to $q\in\mathcal P$, and every $d,e,\ell$ sum below is restricted to $\mathcal D$, including the inner sum in the minimizing weights:
\[
 G(z)=\sum_{\substack{d\in\mathcal D\\d\le z}}H(d),\qquad
 \lambda_d=\frac{\mu(d)}{G(z)}
 \prod_{q\mid d}(1-g(q))^{-1}
 \sum_{\substack{\ell\in\mathcal D,\ \ell\le z/d\\(\ell,d)=1}}H(\ell)
 \quad(d\in\mathcal D,d\le z),
\]
and $\lambda_d=0$ otherwise. Before any CRT expansion, the square majorant is
\[
 \#\{n\in I:n\text{ survives}\}
 \le\sum_{n\in I}
 \left(\sum_{\substack{d\in\mathcal D\\d\le z}}
             \lambda_d\,1_{\Omega_d}(n)\right)^2.
\]
Indeed on a surviving $n$ only the $d=1$ indicator remains and $\lambda_1=1$. Only after this inequality is established does one expand it into the double sum over $d,e\in\mathcal D$ and count $\Omega_{[d,e]}$ by CRT. All subsequent diagonalizing sums over $\ell$ have the same restriction. This is a notation and order-of-proof repair, not an additional pointwise hypothesis.

\section{Executable scope and nonclaims}
Python 3.9 or later and its standard library suffice. The principal commands are
\begin{lstlisting}
python3 verifier.py --self-test --anchors --cover-bound 10000000000
python3 verifier.py --witness 11259889
python3 verifier.py --closure-prime 2521
\end{lstlisting}
The first regenerates all finite angular exceptions, the hand tables, map/inverse tests, the three exact finite-family failures, and the compressed covering verification. The second reconstructs a complete raw ordered witness with a rule identifier. The third exhausts the specified infinite closure and four middle shells for one prime; failure of a stored covering rule alone is never treated as nonexistence.

The cover file is compressed JSON, not executable code. Integers can exceed $2^{53}$; use arbitrary-precision JSON integer parsing. The verifier uses explicit checks rather than optimization-sensitive Python assertions. Primality is certified by complete trial division or a complete sieve, never by a probable-prime test. The Lean-ready plan supplied separately is an exact dependency specification, not a compiled Lean proof. No conclusion about primes above $10^{10}$ is drawn from the census; individual stored rules remain valid wherever their explicit intervals and congruences hold.

\begin{thebibliography}{9}
\bibitem{Yamamoto}
Yamamoto, K. (1965). On the Diophantine equation $4/n=1/x+1/y+1/z$.
\emph{Memoirs of the Faculty of Science, Kyushu University, Series A, Mathematics, 19}(1), 37--47.
\url{https://www.jstage.jst.go.jp/article/kyushumfs/19/1/19_1_37/_article}.

\bibitem{MV}
Monks, M., \& Velingker, A. (2008, September 23).
\emph{On the Erd\H{o}s--Straus conjecture: Properties of solutions to its underlying Diophantine equation} [Author-hosted manuscript].
\url{https://mathematicalgemstones.com/maria/papers/ErdosStraus.pdf}.

\bibitem{ET}
Elsholtz, C., \& Tao, T. (2013). Counting the number of solutions to the Erd\H{o}s--Straus equation on unit fractions.
\emph{Journal of the Australian Mathematical Society, 94}(1), 50--105.
\url{https://arxiv.org/abs/1107.1010}.

\bibitem{Dahan}
Dahan, B. (2026, August 25).
\emph{Sieve dimension and search depth for the Erd\H{o}s--Straus conjecture, $n\equiv1\pmod{24}$} [Preprint, arXiv:2608.24035v1].
\url{https://arxiv.org/html/2608.24035v1}.

\bibitem{Repo11}
chasebryan. (2026, August 15).
\emph{Exact $q=11$ filter for Mordell-hard primes in the strong/Type-II corridor} [Repository research note, \texttt{centl}].
\url{https://github.com/chasebryan/centl/blob/main/research/erdos-straus/STRONG-ES-Q11-EXACT-FILTER.md}.

\bibitem{Repo23}
chasebryan. (2026, August 15).
\emph{Exact $q=23$ filter for Mordell-hard primes in the strong/Type-II corridor} [Repository research note, \texttt{centl}].
\url{https://github.com/chasebryan/centl/blob/main/research/erdos-straus/STRONG-ES-Q23-EXACT-FILTER.md}.
\end{thebibliography}
\end{document}
